\chapter*{Abstract}

Humans explore their visual world through a rapid sequence of eye movements, known as scanpaths, to selectively process information. While deep learning has significantly advanced our ability to predict these gaze patterns, state-of-the-art models are fundamentally limited by their reliance on an implicit, feature-based understanding of scene content. 
They lack an explicit representation of the scene's compositional structure—the very objects and surfaces that are the primary units of human attention.

This thesis posits that a model's predictive power can be substantially improved by explicitly injecting information about a scene's visual segmentation.
To test this hypothesis, we develop a new family of probabilistic models for scanpath prediction, culminating in \textbf{DinoGaze-SPADE}.
This novel architecture introduces two key innovations.
First, it replaces the conventional CNN backbone with a DINOv2 Vision Transformer, leveraging its powerful, self-supervised features that exhibit an emergent, implicit knowledge of scene structure. Second, and most critically, we introduce a novel mechanism called \textbf{"semantic painting"} to explicitly inject structural information. This technique uses a given unsupervised segmentation mask to create a dynamic, feature-rich guidance map that modulates the network's internal feature processing via Spatially-Adaptive Normalization (SPADE), solving the "semantic incoherence" problem that prevents standard methods from working with unsupervised masks.

We rigorously evaluate our models on the standard MIT1003 and SALICON benchmarks.
Our results demonstrate that the ViT-powered backbone alone (DinoGaze) establishes a new, powerful state-of-the-art baseline, outperforming the previous best model, DeepGaze III.
Furthermore, we show that the explicit injection of segmentation information via our semantic painting mechanism provides a consistent and significant additional performance boost, particularly on the complex, sequential task of scanpath prediction. The final DinoGaze-SPADE model represents the new state of the art and, more importantly, provides a validated framework for future research into personalized gaze prediction by enabling the injection of subject-specific perceptual segmentation maps.
This work provides strong computational evidence that an explicit understanding of scene structure is a critical component for accurately and robustly modeling human visual attention.
\thispagestyle{empty}
\mbox{}
\newpage
